\section{Introduction}
\label{sec:intro}

An exposure time calculator answers two reciprocal questions: what
signal-to-noise ratio (S/N) a given exposure of a given target will reach, and
how long an exposure is needed to reach a requested S/N. \spetc{} answers both
for broad-band aperture photometry and for slit or slitless spectroscopy, for
any telescope--detector combination the user describes, at any site, time and
sky position.

The code descends from the Fortran \code{interpola\_spad} programs (2007) and
retains their template-calibration convention
(Sect.~\ref{sec:template}), while the physics has been rebuilt in Python on
\code{astropy} \citep{astropy2022}: all radiometric quantities carry explicit
units, all coordinate and time transformations use IERS-based ERFA routines,
and every empirical relation is taken from the published literature cited in
the section where it is used.

Version 10 upgrades the physical model in six areas, each described in its own
section below:
\begin{enumerate}
  \item a position-dependent night-sky model with van Rhijn airglow, zodiacal
        light, integrated starlight, and the \citet{krisciunas1991} moonlight
        model with the published scattering function (Sect.~\ref{sec:sky});
  \item atmospheric refraction and the \citet{pickering2002} airmass
        (Sect.~\ref{sec:airmass});
  \item selectable Gaussian or \citet{moffat1969} point-spread functions
        (Sect.~\ref{sec:psf});
  \item a noise budget extended with \citet{young1967} scintillation,
        \citet{filippenko1982} atmospheric dispersion and ADC quantization
        noise (Sects.~\ref{sec:dispersion}, \ref{sec:noise});
  \item transmission-grating (Star Analyser 100/200) slitless spectroscopy
        (Sect.~\ref{sec:slitless});
  \item extended sources: galaxies, nebulae and planets treated as uniform
        surface brightness (Sect.~\ref{sec:extended}).
\end{enumerate}

Throughout, wavelengths are in \AA{}ngstr\"om internally, fluxes are
$\flam$ in \si{erg.s^{-1}.cm^{-2}.\angstrom^{-1}}, and detected signals are in
photo-electrons. The notation is summarized in Table~\ref{tab:symbols}.

\begin{table}
\centering
\caption{Principal symbols.}
\label{tab:symbols}
\begin{tabular}{ll}
\toprule
Symbol & Meaning \\
\midrule
$\flam$        & spectral flux density [\si{erg.s^{-1}.cm^{-2}.\angstrom^{-1}}] \\
$A$            & unobstructed collecting area $\pi(D^2-d^2)/4$ \\
$T_\lambda$    & filter transmission (0--1) \\
$Q_\lambda$    & detector quantum efficiency (0--1) \\
$\eta$         & scalar optics throughput excluding QE \\
$O_\lambda$    & optional calibrated instrument response curve \\
$T_{\rm atm}(\lambda)$ & zenith atmospheric transmission \\
$X$            & airmass \\
$z$            & zenith distance; $h$ apparent altitude \\
$s$            & seeing FWHM [arcsec] \\
$q$            & parallactic angle \\
$S,B,D$        & source, sky, dark electrons in the aperture \\
$\sigma_R$     & read noise per pixel [e$^-$ rms] \\
$g$            & gain [e$^-$/ADU] \\
$N_{\rm pix}$  & pixels in the photometric aperture or extraction \\
$R$            & resolving power $\lambda/\Delta\lambda$ \\
\bottomrule
\end{tabular}
\end{table}
